\subsection{Morfismos regulares, normales, reducidos, suaves.}
\label{subsection:IV.6.8}

\begin{definition}[6.8.1]
\label{IV.2.s6.8.1-fr}
Sean $X$, $Y$ dos preesquemas, $f:X\to Y$ un morfismo tal que la fibra
$f^{-1}(y)$ sea un preesquema localmente noetheriano para todo $y\in Y$, $x$
un punto de $X$. Se dice respectivamente que $f$ es un morfismo:
\begin{enumerate}
\item[\textup{(i)}] de coprofundidad $\leq n$ en el punto $x$;
\item[\textup{(ii)}] de Cohen-Macaulay en el punto $x$;
\oldpage[IV-2]{151}
\item[\textup{(iii)}] $(S_n)$ en el punto $x$;
\item[\textup{(iv)}] regular en el punto $x$;
\item[\textup{(v)}] $(R_n)$ en el punto $x$;
\item[\textup{(vi)}] normal en el punto $x$;
\item[\textup{(vii)}] reducido en el punto $x$;
\end{enumerate}
si $f$ es plano en el punto $x$, y si además la propiedad correspondiente
$\boldsymbol{P}(f^{-1}(f(x)),x)$ (notación de \textup{(6.7.8)}) es verdadera.

Se dice que $f$ es suave en el punto $x$ si es localmente de presentación finita
en una vecindad de $x$ en $X$, y si es regular en el punto $x$. Se dice
respectivamente que $f$ es un morfismo: de coprofundidad $\leq n$, de
Cohen-Macaulay, $(S_n)$, regular, $(R_n)$, normal, reducido, suave, si posee la
propiedad correspondiente en todo punto de $X$.
\end{definition}
